\section{Introduction}
A benchmark score is interpretable only together with the procedure that produced the evaluated system. If a search evaluates many prompts on the same questions and returns the empirical winner, the winning score reflects both performance and favorable sampling noise. Reporting that score as if the prompt had been fixed before evaluation can distort uncertainty and subsequent comparisons. The relevant question is quantitative: how much distortion occurs under a specified search, and what resources are needed for a more reliable report?

We study exhaustive automated selection from a frozen bank of configurations. This restricted setting separates selection from candidate generation: the bank exists before evaluation outcomes are observed. It is an authentic search procedure, but it does not represent every compiler that constructs instructions, demonstrations, or pipelines using feedback. In particular, our fixed-bank guarantees do not extend merely by counting the prompts produced by an adaptive generator.

Our central target is the fresh-question accuracy of the configuration actually chosen by the search. This differs from the best available configuration's accuracy and from average quality over repeated searches. These distinctions determine whether an interval, a holdout estimate, and a claim of improvement answer the same question. A reporting correction can improve calibration without improving the chosen system, while reserving evaluation questions can change which system is chosen.

The underlying selection effect and its standard remedies have established statistical foundations. Our contribution is an empirical assessment that connects their assumptions, calibration, decision consequences, and resource costs in a precisely specified configuration-search setting. Controlled experiments isolate sample size, bank size, shared errors, and quality differences. We compare reused-score reporting, simultaneous fixed-bank intervals, and independent auditing, including splits that preserve either total questions or total model calls. We distinguish their targets and report uncertainty from repeated whole-search experiments.

Our real-model study uses newly generated responses on MMLU-Pro and separates three kinds of evidence: a disjoint audit for the selected system on the remaining benchmark, conditional reselection on one response matrix, and reporting calibration against known empirical-cache truth. They support different claims. The real setting exhibits a substantial reuse--audit gap and an audit-supported improvement over default, but does not reproduce the synthetic large-bank undercoverage. This combination motivates reporting the setting and target rather than assuming that every optimized score has the same failure.
